% ============================================================
% IV. RESULTADOS
% ============================================================

\section{Resultados}

\subsection{Item 1 — Muestra01.csv}

La Fig.~\ref{fig:muestra01} muestra la señal completa en el tiempo y el espectro de cada uno de los cuatro segmentos de $0.05$ s. Las frecuencias dominantes obtenidas fueron:

\begin{table}[!t]
\centering
\caption{Frecuencia dominante por segmento de Muestra01.csv}
\label{tab:muestra01}
\small
\begin{tabular}{@{}l r@{}}
\toprule
Segmento & $f_d$ (Hz) \\
\midrule
1 & 1200.00 \\
2 & 1600.00 \\
3 & 940.00 \\
4 & 1300.00 \\
\bottomrule
\end{tabular}
\end{table}

\begin{figure}[H]
\centering
\includegraphics[width=\linewidth]{muestra01_analisis.png}
\caption{Señal completa de Muestra01.csv en el tiempo y espectro de cada uno de sus cuatro segmentos de $0.05$ s.}
\label{fig:muestra01}
\end{figure}

Las cuatro frecuencias dominantes son distintas entre sí, lo que sugiere que la señal registrada no es estacionaria: su contenido espectral dominante cambia de un segmento de $0.05$ s al siguiente.

\subsection{Item 2 — Muestreo, cuantización y aliasing}

La Fig.~\ref{fig:muestreo_cuantizacion} muestra la señal seno original de $100$ Hz (línea punteada) junto con las muestras cuantizadas a $3$ bits para cada frecuencia de muestreo evaluada.

\begin{table}[!t]
\centering
\caption{Resultado del muestreo del seno de 100 Hz}
\label{tab:muestreo}
\small
\begin{tabular}{@{}r r l@{}}
\toprule
$f_s$ (Hz) & Muestras (0.05 s) & Condición de Nyquist \\
\midrule
70 & 4 & $f_s < 200$ Hz — \textbf{aliasing} \\
500 & 25 & $f_s \geq 200$ Hz — sin aliasing \\
1000 & 50 & $f_s \geq 200$ Hz — sin aliasing \\
\bottomrule
\end{tabular}
\end{table}

\begin{figure}[H]
\centering
\includegraphics[width=\linewidth]{muestreo_cuantizacion.png}
\caption{Señal seno de 100 Hz muestreada y cuantizada (3 bits) a 70, 500 y 1000 Hz. A 70 Hz las muestras son demasiado escasas para seguir la forma de onda real: el patrón muestreado es consistente con una frecuencia mucho menor a 100 Hz (aliasing).}
\label{fig:muestreo_cuantizacion}
\end{figure}

Con solo $4$ muestras en $0.05$ s, el muestreo a $70$ Hz no captura ni siquiera dos ciclos completos de la señal de $100$ Hz, y las muestras resultantes son compatibles con una señal de frecuencia mucho menor —el efecto de aliasing descrito en la Sección~\ref{sec:aliasing}. A $500$ y $1000$ Hz, con $5$ y $10$ muestras por ciclo respectivamente, la forma de onda queda claramente representada.

\subsection{Item 3 — FFT de violín, tambor y gato}

\begin{table}[!t]
\centering
\caption{Resultados del análisis espectral (señal completa)}
\label{tab:audio}
\small
\begin{tabular}{@{}l r r r@{}}
\toprule
Señal & $f_s$ (Hz) & Duración (s) & $f_d$ (Hz) \\
\midrule
Violín & 44100 & 2.00 & 385.50 \\
Tambor & 44100 & 0.96 & 7596.10 \\
Gato & 11025 & 1.58 & 1715.08 \\
\bottomrule
\end{tabular}
\end{table}

\begin{figure}[H]
\centering
\includegraphics[width=\linewidth]{audio_fft_comparacion.png}
\caption{Espectros de magnitud (FFT) del violín, el tambor y el gato, superpuestos.}
\label{fig:audio_fft}
\end{figure}

Las tres frecuencias dominantes están claramente separadas entre sí: el violín concentra su energía en $385.50$ Hz (rango de una nota musical grave/media, coherente con la nota fundamental de una cuerda frotada), el gato en $1715.08$ Hz (coherente con el tono agudo de un maullido), y el tambor —de forma notable— en $7596.10$ Hz, muy por encima de lo que se esperaría de la frecuencia fundamental de un parche de tambor. Esto se debe a que el ataque de un tambor es un transitorio de banda ancha con energía significativa en armónicos altos y en el ruido del golpe, por lo que el pico de \emph{mayor magnitud} de su espectro no corresponde necesariamente a un tono fundamental grave perceptible, sino al contenido de alta frecuencia del transitorio de ataque —un punto que se retoma en la Discusión.
